Model selection is divided between two cooperating agents, coordinated by an orchestrator that routes work, retries failed steps, and determines when to stop. The Retrieval Agent identifies promising candidates, while the Evaluation Agent verifies their performance on the target data.

\paragraph{Retrieval Agent: making a cold-start dataset scorable.}
A cold-start dataset has no existing node or evaluation edges in the artifact graph, and therefore cannot be scored directly. To simulate this setting, we use a new Hugging Face dataset that is absent from ArtifactLinker's graph. Both retrieval strategies share two preprocessing steps. First, GPT-4o \citep{openai2024gpt4o} summarizes the dataset card into a structured description of its task, domain, size, and format, using the same prompt employed by ArtifactLinker to summarize its dataset nodes. Second, we embed this summary with Voyage AI's \texttt{voyage-3} model \citep{voyageai2024voyage3}, matching ArtifactLinker's node features and placing the target dataset in the same representation space as the 1{,}208 existing dataset nodes.
The strategies then diverge in how they use that graph:

\textbf{Approach 1 -- cosine retrieval.} We compute cosine similarity between the target dataset and all existing dataset nodes, retrieving the most similar datasets and their associated model evaluations. These models, together with their measured scores and metrics, form a \emph{graph-conditioned} prompt that GPT-5.5 \citep{openai2026gpt55} uses to reason across similar datasets and recommend the top-5 candidate models, with a justification for each.

\textbf{Approach 2 -- direct GNN scoring.} We insert the target dataset into ArtifactLinker's graph as a new node, using its \texttt{voyage-3} embedding as the node feature, and apply the trained GNN to rank candidate models by their predicted link probability $\times$ evaluation score. This tests ArtifactLinker inductively, as the target dataset was unseen during training. The top-ranked models and their supporting evaluation evidence are then passed to GPT-5.5 \citep{openai2026gpt55} to recommend the top-5 candidate models, with a justification for each.

Both approaches are constraint-aware: candidates that violate the user's stated constraints (e.g., parameter budget) are filtered before ranking. The retrieval LLM receives two distinct evidence blocks -- ArtifactLinker (GNN) context and dataset-similarity context -- as separate, labeled sources, so its recommendation can be traced back to which evidence supported which pick.

\paragraph{Why both?}
The two approaches fail in different regions of the graph. Cosine retrieval is only as good as the target's nearest neighbors: when the graph contains closely related datasets (e.g., a new sentiment or news-classification set next to many well-evaluated ones), it hands the LLM directly measured scores on near-identical tasks -- interpretable, metric-aware evidence that is hard to beat. But it sees only the models evaluated on those few neighbors, and when the target lies in a sparse region (a niche medical or legal task whose closest neighbor is unrelated), that evidence is thin or misleading. Direct GNN scoring instead ranks every model in the graph, aggregating multi-hop signal -- a model's results across many datasets, its fine-tuning lineage, shared papers and code -- so it can credit models that transfer broadly even when no single neighbor is close. Its weaknesses are the mirror image: the inductive node carries only its text embedding, its scores are opaque regression outputs rather than measured metrics, and the link head favors heavily evaluated, popular models regardless of fit. We therefore expect cosine retrieval to win where the neighborhood is dense and GNN scoring where it is sparse or where the best model's evidence is structurally rather than semantically adjacent; the constraint-aware merge lets the LLM weigh both, and the ablations in Section~\ref{sec:experimental_setup} isolate each signal's contribution.

\paragraph{Evaluation Agent: measuring performance de facto.}
The Evaluation Agent takes the candidate models and the user's dataset, plans an evaluation, writes and runs the evaluation code, handles execution errors and retries, and records the measured score for each candidate -- the same verify step ArtifactLinker applies to its own top-ranked links, now applied to our retrieval agent's cold-start shortlist.

\paragraph{Closed-loop refinement.}
Model selection proceeds in two runs. In the \emph{initial run}, the Retrieval Agent produces a ranked shortlist and the Evaluation Agent measures each candidate on the user's data. The orchestrator then analyzes where each candidate fails (which classes, slices, or examples it gets wrong) and launches a \emph{revision run}: the Retrieval Agent is re-prompted with the original graph evidence, the initial-run shortlist, and the measured results with error notes, and confirms, drops, or replaces candidates based on what actually happened rather than re-predicting from the graph alone. Selection thus becomes iterative experimentation that targets its own failures instead of a single blind pass.
